\section{反向传播算法}

\subsection{反向传播的本质}

\begin{importantbox}{反向传播算法 = 两件事}
\begin{enumerate}
    \item \textbf{使用链式法则}：将复合函数的梯度分解为局部梯度的乘积
    \item \textbf{存储中间结果}：避免重复计算——同一个上游梯度被多个下游节点共享
\end{enumerate}
这就是反向传播算法的全部。
\end{importantbox}

\subsection{计算图}

为了系统化地实现反向传播，我们将计算过程表示为\textbf{计算图}（Computation Graph）：

\begin{figure}[H]
\centering
\includegraphics[width=0.95\textwidth]{slides-images/slide-050.jpg}
\caption{NER 网络的计算图：源节点是输入（$x, W, b, u$），内部节点是操作（$\cdot, +, f, \cdot$），边传递中间结果。前向传播沿图从左到右计算函数值\protect\footnotemark}
\end{figure}
\footnotetext{Slides 第51页。}

计算图中：
\begin{itemize}
    \item \textbf{源节点}：输入变量（$x, W, b, u$）
    \item \textbf{内部节点}：运算操作（矩阵乘法、加法、激活函数、点积）
    \item \textbf{边}：传递操作的结果
\end{itemize}

\textbf{前向传播}（Forward Propagation）：沿计算图从左到右，依次计算每个节点的值。

\textbf{反向传播}（Backward Propagation / Backpropagation）：沿计算图从右到左，依次计算梯度。

\begin{figure}[H]
\centering
\includegraphics[width=0.95\textwidth]{slides-images/slide-052.jpg}
\caption{反向传播（Backward Pass）：梯度沿图从右到左传播\protect\footnotemark}
\end{figure}
\footnotetext{Slides 第53页。}

\subsection{单节点的梯度传播规则}

对于计算图中的任意一个节点，设其计算为 $h = f(z)$：

\begin{figure}[H]
\centering
\includegraphics[width=0.95\textwidth]{slides-images/slide-054.jpg}
\caption{单节点梯度传播：下游梯度 = 上游梯度 $\times$ 局部梯度\protect\footnotemark}
\end{figure}
\footnotetext{Slides 第55页。}

$$\underbrace{\frac{\partial s}{\partial z}}_{\text{下游梯度}} = \underbrace{\frac{\partial s}{\partial h}}_{\text{上游梯度}} \times \underbrace{\frac{\partial h}{\partial z}}_{\text{局部梯度}}$$

这就是链式法则在计算图中的体现。当节点有多个输入时，对每个输入分别计算局部梯度，然后分别与上游梯度相乘。

\subsection{实例：$f(x,y,z) = (x+y) \cdot \max(y,z)$}

Manning 教授用一个具体的非神经网络的例子来演示反向传播的过程。

\begin{figure}[H]
\centering
\includegraphics[width=0.95\textwidth]{slides-images/slide-060.jpg}
\caption{计算图实例：$f(x,y,z) = (x+y) \cdot \max(y,z)$，当 $x=1, y=2, z=0$ 时\protect\footnotemark}
\end{figure}
\footnotetext{Slides 第61页。}

\textbf{前向传播}：
\begin{itemize}
    \item $a = x + y = 1 + 2 = 3$
    \item $b = \max(y, z) = \max(2, 0) = 2$
    \item $f = a \cdot b = 3 \times 2 = 6$
\end{itemize}

\textbf{局部梯度}：
\begin{itemize}
    \item 加法节点：$\frac{\partial a}{\partial x} = 1$, $\frac{\partial a}{\partial y} = 1$
    \item max 节点：$\frac{\partial b}{\partial y} = 1$（$y > z$），$\frac{\partial b}{\partial z} = 0$
    \item 乘法节点：$\frac{\partial f}{\partial a} = b = 2$，$\frac{\partial f}{\partial b} = a = 3$
\end{itemize}

\textbf{反向传播}（从右到左）：

\begin{figure}[H]
\centering
\includegraphics[width=0.95\textwidth]{slides-images/slide-063.jpg}
\caption{反向传播过程：$\frac{\partial f}{\partial f} = 1$，乘法节点分发梯度 2 和 3，加法节点分发梯度 2，max 路由梯度 3 到 $y$\protect\footnotemark}
\end{figure}
\footnotetext{Slides 第64页。}

\begin{enumerate}
    \item $\frac{\partial f}{\partial f} = 1$
    \item 乘法节点：$\frac{\partial f}{\partial a} = b = 2$，$\frac{\partial f}{\partial b} = a = 3$
    \item 加法节点：$\frac{\partial f}{\partial x} = \frac{\partial f}{\partial a} \cdot 1 = 2$，从加法到 $y$：$\frac{\partial f}{\partial a} \cdot 1 = 2$
    \item max 节点：到 $y$：$\frac{\partial f}{\partial b} \cdot 1 = 3$，到 $z$：$\frac{\partial f}{\partial b} \cdot 0 = 0$
    \item \textbf{$y$ 有两条路径}：$\frac{\partial f}{\partial y} = 2 + 3 = 5$
\end{enumerate}

最终结果：$\frac{\partial f}{\partial x} = 2$，$\frac{\partial f}{\partial y} = 5$，$\frac{\partial f}{\partial z} = 0$。

\textbf{验证}：将 $y$ 从 2 改为 2.1，则 $a = 3.1$，$b = 2.1$，$f = 3.1 \times 2.1 = 6.51$，变化量 $\approx 0.51 \approx 5 \times 0.1$，与梯度 5 吻合。

\subsection{分支处的梯度求和}

\begin{figure}[H]
\centering
\includegraphics[width=0.95\textwidth]{slides-images/slide-070.jpg}
\caption{分支处梯度求和：当一个变量出现在多个路径中时，各路径的梯度需要\textbf{相加}\protect\footnotemark}
\end{figure}
\footnotetext{Slides 第71页。}

上例中 $y$ 同时参与了加法和 max 两个计算，因此它的总梯度是两条路径上梯度的\textbf{和}：

$$\frac{\partial f}{\partial y} = \frac{\partial f}{\partial a} \cdot \frac{\partial a}{\partial y} + \frac{\partial f}{\partial b} \cdot \frac{\partial b}{\partial y} = 2 \times 1 + 3 \times 1 = 5$$

这就是多元链式法则的体现。

\subsection{三种基本操作的梯度直觉}

\begin{figure}[H]
\centering
\includegraphics[width=0.95\textwidth]{slides-images/slide-072.jpg}
\caption{三种基本操作的梯度行为：加法分发梯度，max 路由梯度，乘法切换梯度\protect\footnotemark}
\end{figure}
\footnotetext{Slides 第73页。}

\begin{importantbox}{三种基本操作的梯度行为}
\begin{itemize}
    \item \textbf{加法}（$+$）：\textbf{分发}梯度——上游梯度原封不动地传给两个输入
    \item \textbf{Max}：\textbf{路由}梯度——上游梯度全部传给较大的那个输入，另一个输入得到零
    \item \textbf{乘法}（$\times$）：\textbf{切换}梯度——上游梯度乘以\textbf{另一个}输入的值传给当前输入
\end{itemize}
\end{importantbox}

\begin{warningbox}{Max 操作导致梯度消失}
Max 操作对``输的那一方''梯度为零——这与 ReLU 的行为一致（$\text{ReLU}(x) = \max(0, x)$）。当 $x < 0$ 时，ReLU 的梯度为零，意味着该方向的信息完全被阻断。这从计算图的角度解释了 dying ReLU 现象。
\end{warningbox}

\subsection{高效反向传播：避免重复计算}

\begin{figure}[H]
\centering
\includegraphics[width=0.95\textwidth]{slides-images/slide-074.jpg}
\caption{低效 vs 高效的梯度计算：逐个计算 $\frac{\partial s}{\partial b}, \frac{\partial s}{\partial W}, \ldots$ 会产生大量重复计算（蓝色/红色重叠部分）\protect\footnotemark}
\end{figure}
\footnotetext{Slides 第75页。}

如果分别独立地计算 $\frac{\partial s}{\partial b}$、$\frac{\partial s}{\partial W}$、$\frac{\partial s}{\partial x}$、$\frac{\partial s}{\partial u}$，会发现它们共享大量的中间计算（上游梯度 $\delta$）。反向传播的核心效率在于：\textbf{按照拓扑排序的逆序一次性遍历计算图}，每个中间梯度只计算一次。

\begin{importantbox}{反向传播的计算复杂度}
如果正确实现，反向传播的时间复杂度与前向传播的时间复杂度\textbf{同阶}（Big-O 相同）。如果你的反向传播比前向传播慢很多，说明你在某处做了重复计算。
\end{importantbox}

\subsection{通用反向传播算法}

\begin{figure}[H]
\centering
\includegraphics[width=0.95\textwidth]{slides-images/slide-076.jpg}
\caption{通用反向传播算法：前向传播按拓扑排序计算函数值，反向传播按逆拓扑序计算梯度\protect\footnotemark}
\end{figure}
\footnotetext{Slides 第77页。}

算法流程：

\textbf{前向传播}：
\begin{enumerate}
    \item 对计算图做\textbf{拓扑排序}，确保每个节点只依赖已计算的节点
    \item 按拓扑序遍历，对每个节点调用其 \texttt{forward} 方法计算输出值
\end{enumerate}

\textbf{反向传播}：
\begin{enumerate}
    \item 初始化输出梯度：$\frac{\partial s}{\partial s} = 1$
    \item 按\textbf{逆拓扑序}遍历节点
    \item 对每个节点：$\text{下游梯度} = \text{上游梯度} \times \text{局部梯度}$
    \item 对分支节点：将多条路径的梯度\textbf{求和}
\end{enumerate}

该算法适用于任意有向无环图（DAG），不要求网络是规整的层状结构。

\subsection{本章小结}

反向传播算法是链式法则在计算图上的高效实现。其核心是``一次前向、一次反向''，利用中间结果共享避免重复计算。对于任意 DAG 结构的计算图，反向传播的复杂度与前向传播同阶。

%% ========== Part 7: 自动微分与实现 ==========

